\section{Asteroid Collision}
\label{sec:sq-asteroid-collision}

\problemstatement{We are given an array $\textit{asteroids}$, where each
element's absolute value is an asteroid's size and its sign is its
direction of travel (positive moves right, negative moves left; all move
at the same speed). When two asteroids meet, the smaller one explodes; if
both are the same size, both explode. Asteroids moving in the same
direction never meet. Return the state of the asteroids after all
collisions have finished.}

\begin{patternbox}
\emph{``Process a sequence left to right, where a new element can
repeatedly `defeat' and remove previously-kept elements (or be defeated
by them) based on a comparison''} is, structurally, the same skeleton as
next greater element (Section~\ref{sec:sq-next-greater-1}): a stack of
``currently surviving'' items, where each new arrival potentially knocks
out some suffix of what's already on the stack before either joining it
or being knocked out itself.
\end{patternbox}

\subsection*{Understanding the problem carefully}

A collision can only happen between a \textbf{right-moving} asteroid that
was placed earlier and a \textbf{left-moving} asteroid that arrives later
(two right-movers never catch up to each other since they move at the
same speed; a left-mover followed by a right-mover are moving apart and
never meet). So, scanning left to right and keeping a stack of asteroids
that have survived so far: whenever a new \emph{left-moving} asteroid
arrives and the stack's top is a surviving \emph{right-moving} asteroid,
they are on a collision course, and we must resolve who survives --
possibly triggering a chain of further collisions if the new asteroid is
large enough to destroy multiple right-movers in a row.

\begin{ideabox}[Stack of Surviving Asteroids, with Chain Collisions]
Maintain a stack. For a right-moving asteroid (positive), simply push it
-- no immediate collision is possible yet. For a left-moving asteroid
(negative): while the stack's top is positive (a right-mover) and smaller
in size than the incoming asteroid, pop it (it loses). If the stack's
top is positive and exactly equal in size, pop it and the incoming
asteroid is also destroyed (mutual annihilation) -- stop processing this
asteroid. If the stack's top is positive and larger, the incoming asteroid
is destroyed -- stop processing this asteroid, do not push it. If the
stack becomes empty or its top is negative (no more right-movers to
collide with), and the incoming asteroid was never destroyed, push it.
\end{ideabox}

\subsection*{Why processing collisions as a chain (not just one comparison) is correct}

A single incoming left-moving asteroid might be large enough to destroy
\emph{several} consecutive right-moving asteroids already on the stack
before either finally being destroyed itself or surviving to be pushed.
Each such destruction is resolved by exactly the comparison rule above,
applied repeatedly -- this is precisely the same ``pop everything the new
arrival invalidates'' loop structure as a monotonic stack, except the
``invalidation'' rule here is a size comparison between asteroids of
opposite direction rather than a simple value comparison, and the loop
must also account for the three-way outcome (incoming destroyed, top
destroyed, or both destroyed) rather than a single binary pop/no-pop
decision.

\subsection*{Algorithm}

\begin{enumerate}
  \item Initialize an empty stack.
  \item For each asteroid $a$ in $\textit{asteroids}$:
  \begin{enumerate}
    \item $\textit{alive} \gets \text{true}$.
    \item While $\textit{alive}$, $a < 0$, the stack is not empty, and
          the stack's top is $>0$:
    \begin{itemize}
      \item If top $< -a$: pop (the top asteroid is destroyed); continue
            the loop.
      \item Else if top $= -a$: pop; $\textit{alive} \gets \text{false}$
            (mutual destruction).
      \item Else (top $> -a$): $\textit{alive} \gets \text{false}$
            (the incoming asteroid is destroyed).
    \end{itemize}
    \item If $\textit{alive}$: push $a$.
  \end{enumerate}
  \item Return the stack's contents, bottom to top.
\end{enumerate}

\subsection*{Dry run}

\dryrun{Take $\textit{asteroids} = [5, 10, -5]$.}

\begin{itemize}
  \item $a=5$: positive, push. Stack $=[5]$.
  \item $a=10$: positive, push. Stack $=[5,10]$.
  \item $a=-5$: negative. Top is $10>0$. Is $10 < 5$? No. Is $10=5$?
        No. So top ($10$) $> 5$: the incoming asteroid ($-5$) is
        destroyed. \textit{alive} $=$ false. Do not push.
\end{itemize}

Final stack: $[5,10]$.

Now take $\textit{asteroids} = [8,-8]$: $a=8$ pushed. $a=-8$: top is $8$,
and $8 = 8$: pop, mutual destruction, \textit{alive}$=$false. Final
stack: $[]$ (both annihilated).

\subsection*{C++ implementation}

\begin{lstlisting}
#include <bits/stdc++.h>
using namespace std;

vector<int> asteroidCollision(vector<int>& asteroids) {
    vector<int> st; // used as a stack via back()/push_back()/pop_back()

    for (int a : asteroids) {
        bool alive = true;
        while (alive && a < 0 && !st.empty() && st.back() > 0) {
            if (st.back() < -a) {
                st.pop_back(); // top destroyed, continue checking
            } else if (st.back() == -a) {
                st.pop_back();
                alive = false; // mutual destruction
            } else {
                alive = false; // incoming asteroid destroyed
            }
        }
        if (alive) st.push_back(a);
    }
    return st;
}

int main() {
    vector<int> asteroids = {5, 10, -5};
    for (int x : asteroidCollision(asteroids)) cout << x << " ";
    cout << endl;
    return 0;
}
\end{lstlisting}

\begin{complexitybox}
\textbf{Time Complexity.} Each asteroid is pushed at most once and popped
at most once across the entire scan, giving
\[
  O(n).
\]
\textbf{Space Complexity.} $O(n)$ for the stack in the worst case (no
collisions occur).
\end{complexitybox}

\begin{takeawaybox}
Monotonic-stack-style ``pop everything the new arrival invalidates''
logic generalizes beyond simple value comparisons: here the invalidation
rule is a three-way size comparison between opposite-direction asteroids,
but the overall shape -- a while-loop popping a suffix of the stack, then
conditionally pushing the new arrival -- is unchanged from
Section~\ref{sec:sq-next-greater-1}.
\end{takeawaybox}
